\chapter{Area of Responsibility}

In vFIR Warszawa, Military Air Traffic Control service can be provided only by vATCOs holding a valid military endorsement.

The following airports are considered military airports:
\begin{itemize}
    \item Cewice (EPCE),
    \item Darłowo (EPDA),
    \item Dęblin (EPDE),
    \item Inowrocław (EPIR),
    \item Krzesiny (EPKS),
    \item Łask (EPLK),
    \item Łęczyca (EPLY),
    \item Malbork (EPMB),
    \item Mirosławiec (EPMI),
    \item Mińsk Mazowiecki (EPMM),
    \item Oksywie (EPOK),
    \item Pruszcz Gdański (EPPR),
    \item Powidz (EPPW),
    \item Radom (EPRA)\footnote{only when military approach service is provided},
    \item Świdwin (EPSN),
    \item Tomaszów Mazowiecki (EPTM).
\end{itemize}

Military controlled airspaces (MCTRs, MTMAs) are \textbf{class D} airspaces.

General Air Traffic (GAT) is conducted according to rules presented in chapters~\ref{ch:twr:proc} and~\ref{ch:app:methods}.

\section{Top-down coverage}
In the absence of military ATC, MCTR and MTMA are relegated to uncontrolled airspace (class G). MCTRs and MTMAs are not covered top-down by higher ATC positions.

\section{Departure and arrival at a military airport}
Aircraft conducting a GAT flight should use the appropriate published SID/STAR fixes as entry/exit points or, if these are unavailable, any fix within 50 NM of the ARP.\@ The selected exit fix should allow to join the rest of the route without additional maneuvers.

When using FRA, the entry/exit point is any FRA INTERMEDIATE point within 50 NM of the ARP.\@

Aircraft conducting an OAT flight should use OAT coordination points appropriate for the airport as entry/exit points. The list of OAT coordination points is available in \href{https://www.ais.pansa.pl/mil/pliki/EP_ENR_2_3_0-1_en.pdf}{AIP MIL ENR 2.3.0}.